\subsubsection{Plan de desarrollo}

El plan traduce los requerimientos priorizados en tarjetas de un tablero Kanban, ordenadas por dependencia. Cada tarjeta corresponde, en general, a un módulo (Tabla~\ref{tab:ob2_modulos}) y se cierra con sus pruebas.

La Tabla~\ref{tab:ob2_tarjetas} presenta las quince tarjetas. Las tres tarjetas en negrita corresponden al agente conversacional. A octubre de 2026 están terminadas las seis primeras; las nueve restantes están pendientes.

\begingroup
\fontsize{9}{11}\selectfont
\renewcommand{\arraystretch}{1.2}
\setlength{\LTpre}{0pt}\setlength{\LTpost}{6pt}
\setlength{\tabcolsep}{4pt}
\begin{longtable}{|p{0.24\linewidth}|p{0.17\linewidth}|p{0.17\linewidth}|p{0.10\linewidth}|p{0.09\linewidth}|p{0.10\linewidth}|}
\caption{Tarjetas del tablero Kanban, con sus requerimientos, prioridad, dependencias y estado a octubre de 2026}
\label{tab:ob2_tarjetas} \\
\hline
\textbf{Tarjeta} & \textbf{Módulo} & \textbf{Requerimientos} & \textbf{Prioridad} & \textbf{Depende de} & \textbf{Estado} \\ \hline
\endfirsthead
\hline
\textbf{Tarjeta} & \textbf{Módulo} & \textbf{Requerimientos} & \textbf{Prioridad} & \textbf{Depende de} & \textbf{Estado} \\ \hline
\endhead
\hline
\endfoot

1. Identidad y autenticación & Identidad y acceso & RNF01 & Debe & --- & Terminada \\ \hline
2. Profesionales y pacientes & Pacientes y profesionales & RF06 (parcial), RF07 (base) & Debe & 1 & Terminada \\ \hline
3. Historia clínica y antecedentes & Clínico & RF01 & Debe & 2 & Terminada \\ \hline
4. Agenda y horarios & Agenda & RF07 & Debe & 2 & Terminada \\ \hline
5. Atención clínica y CIE-10 & Clínico & RF01, RF03, RF06 & Debe & 3, 4 & Terminada \\ \hline
6. Odontograma & Clínico & RF02 & Debe & 5 & Terminada \\ \hline
7. Tratamientos & Clínico & RF04 & Debe & 5, 6 & Pendiente \\ \hline
8. Recetas & Clínico & RF04, RF05 & Debe & 5 & Pendiente \\ \hline
9. Cobros y comprobante & Cobros & RF08 & Debe & 7 & Pendiente \\ \hline
10. Auditoría & Transversal & RNF02 & Debe & 9 & Pendiente \\ \hline
11. Sincronización con Google Calendar &Agenda & RF07 & Debe & 4 & Pendiente \\ \hline
\textbf{12. Agente: núcleo, aviso, información y derivación} & Agente & RF09, RF10, RF13, RNF03, RNF05 & Debe y Debería & 4 & Pendiente \\ \hline
\textbf{13. Agente: agendar y reprogramar} & Agente & RF11, RF12 & Debe & 12 & Pendiente \\ \hline
\textbf{14. Agente: recordatorios} & Agente & RF14 & Debería & 13 & Pendiente \\ \hline
15. Respaldos y exportación & Servicio de respaldo & RNF04 & Debe & Todas & Pendiente \\ \hline
\end{longtable}
\endgroup

\paragraph{Estrategia de pruebas}

A la fecha, el backend cuenta con 30 clases de prueba: 13 de integración y 17 unitarias, que suman 118 métodos de prueba. Las pruebas de integración se ejecutan contra una base de datos PostgreSQL 16 en contenedor. La Tabla~\ref{tab:ob2_pruebas} define el alcance y el criterio de cada nivel.

\begin{table}[H]
\centering
\caption{Niveles de prueba del sistema}
\label{tab:ob2_pruebas}
\fontsize{9}{11}\selectfont
\renewcommand{\arraystretch}{1.2}
\setlength{\tabcolsep}{4pt}
\begin{tabular}{|>{\raggedright\arraybackslash}p{0.17\linewidth}|>{\raggedright\arraybackslash}p{0.34\linewidth}|>{\raggedright\arraybackslash}p{0.37\linewidth}|}
\hline
\textbf{Nivel} & \textbf{Alcance} & \textbf{Criterio} \\ \hline
Integración & Endpoints de los módulos críticos (agenda, historia clínica, atención y odontograma) sobre una base de datos real en contenedor & Un flujo principal y un caso de fallo clave por módulo, como el rechazo de una cita en un horario ocupado \\ \hline
Unitaria & Reglas puras: estados de la cita, cálculo de horarios, reglas del agente y máquina de pasos del agente & Se prueban las transiciones permitidas y las rechazadas, sin acceso a la base de datos \\ \hline
Anticruce & Restricción de cruces de citas por profesional, verificada por la API sobre la base de datos & Un intento de cita solapada se rechaza mientras la cita existente no se cancele. La reprogramación pasa por la misma restricción; su prueba sobre un horario ocupado se añade con la tarjeta 13 \\ \hline
Clínica (Objetivo 3) & Recorrido de los flujos con el personal del centro & Los tramos $T_1$ a $T_5$ se comparan con la línea base del pretest \\ \hline
\end{tabular}
\end{table}
